\section{Introduction}
Large Language Models (LLMs) have seen an unprecedented rise in popularity and application across various domains, including education. This integration of LLMs into education contexts has generated considerable interest, driven by their demonstrated proficiency in natural language understanding, text generation, and reasoning tasks \cite{Brown2020}. However, while LLMs excel in natural language processing, one must pose the question: how well can these models generalize beyond language generation to perform effectively in structured, non-linguistic tasks, such as tabular classification? This question carries practical significance for Higher Education Institutions (HEIs), where the ability to predict students' academic outcomes from behavioral and demographic data underpins early-warning systems that identify at-risk students and allocate support resources.

Regarding classification tasks, there exists an extensive, well-established literature on the use of traditional machine learning models, such as Random Forest, Gradient Boosting, and Support Vector Machines, to address such problems. These models are explicitly trained on the target dataset, benefit from hyperparameter tuning, and operate natively on tabular feature vectors. Their effectiveness is well-documented \cite{Delgado2014, Jimenez-Macias2025, Zhan2025, Bujang2021, Alhakbani2022}. LLMs, by contrast, acquire knowledge through large-scale pretraining on textual corpora and are not exposed to structured tabular data during this process. For instance, when presented with a student record serialized as JSON, the LLM must interpret a data format fundamentally different from the natural language it was trained on.

Recent studies have begun examining LLM performance on tabular data. Wen et al. (2025) proposed Tabular In-Context Learning (TabICL), a retrieval-augmented framework that adapts LLMs for tabular classification \cite{Wen2025}. However, such approaches incorporate task-specific mechanisms (retrieval pipelines, structured prompting protocols) that blur the line between general-purpose inference and task-adapted learning. Furthermore, most existing comparisons operate on generic benchmarks (e.g., UCI, OpenML datasets) rather than domain-specific datasets with complex, multi-source feature engineering pipelines, a common characteristic of real-world educational data.

This gap motivates our central research question: To what extent can general-purpose LLMs, without any task-specific adaptation, generalize to a structured classification task, particularly in an educational context? We therefore deliberately exclude fine-tuning techniques (e.g., LoRA, QLoRA) and retrieval-augmented generation (RAG), as these modifications would transform the LLM into a task-adapted model, answering a different research question — namely, how much domain-specific adaptation is required to close the performance gap. To this end, we also favor simple prompting strategies (zero-shot, one-shot, few-shot) that do not require engineering complex prompt templates or retrieval mechanisms, as these would again shift the focus from generalization to adaptation.

To this end, we evaluate two open-weight LLMs, Qwen-2.5-3B and LLaMA-3-8B, using zero-shot, one-shot, and five-shot prompting on the Open University Learning Analytics Dataset (OULAD) \cite{Kuzilek2017} and benchmark these against seven traditional machine learning models that undergo full hyperparameter tuning via cross-validation. We also designed a pipeline to serialize tabular data into a JSON format that is compatible with LLM input, ensuring that the evaluation reflects the LLMs' ability to interpret structured data rather than relying on engineered prompt templates. Our findings reveal that the tested LLMs underperform heavily compared to traditional models, with the best LLM configuration achieving a Macro F1 of 0.481, while the best traditional model reaches 0.722. This significant performance gap highlights the limitations of general-purpose LLMs in structured classification tasks and underscores the need for further research into domain adaptation techniques or architectural modifications to enhance their applicability in non-linguistic contexts.

The remainder of this paper is organized as follows. Section 2 provides background on the models used. Section 3 reviews related work on machine learning classifiers in education, data aggregation, and performance metrics. Section 4 describes the methodology, including the dataset, preprocessing pipeline, and experimental setup. Section 5 presents the results, and Section 6 discusses their implications, limitations, and directions for future work.